%% {{CODE}} manuscript template -- {{PUBNAME}}
%%
%% Submit this file compiled with the `anonymous` option, which removes the
%% author block for double-blind review. Give author names, affiliations and
%% acknowledgements in titlepage.tex instead, and submit that separately.
%% `lineno` adds line numbers for the reviewers; keep it for submission.
\documentclass[journal,anonymous,lineno,{{CITEOPT}}]{nextgeniq}

\publication{{{PUBNAME}}}{{{CODE}}}
\articletype{Research Article} % or Review Article, Conceptual Paper, Case Study, Short Communication
\shorttitle{A short running title of up to 50 characters}

\title{Title of the Article: Concise, Specific and Informative}

% Author details are hidden by the `anonymous` option. Fill them in anyway:
% the same block is used for the accepted version.
\author[1]{First Author}
\author[2]{Second Author}
\author[1]{Third Author}
\affil[1]{Department, University, City, Country}
\affil[2]{Department, Institution, City, Country}
\correspondence{First Author, \texttt{first.author@example.org}}

\keywords{keyword one; keyword two; keyword three; keyword four; keyword five}
% Article history: completed by the editorial office. Leave these as they are.
\received{DD Month YYYY}
\revised{DD Month YYYY}
\accepted{DD Month YYYY}
\published{DD Month YYYY}

\begin{document}

\maketitle

\begin{abstract}
State the problem, the approach, the main findings and their implications in a
single paragraph of 150--250 words. Do not cite references, and spell out
abbreviations at first use. A structured abstract (Purpose, Design/methodology,
Findings, Originality) is also acceptable.
\end{abstract}

\section{Introduction}
\label{sec:intro}

Explain the research question, why it matters, and what the article
contributes. Close the introduction with a short outline of the sections that
follow. Cite sources as \citet{smith2024} or \citep{lee2023,garcia2022}.

\section{Related Work}
\label{sec:related}

Position the study in the existing literature and identify the gap it
addresses \citep{garcia2022}.

\section{Methods}
\label{sec:methods}

Describe the data, materials and procedures in enough detail for the study to
be reproduced. Number displayed equations so they can be referred to, as in
Equation~\eqref{eq:example}.
\begin{equation}
  \hat{y} = \arg\max_{y \in \mathcal{Y}} \; P(y \mid x; \theta)
  \label{eq:example}
\end{equation}

\subsection{Data collection}
Use subsections where they help the reader; avoid going deeper than three
levels.

\section{Results}
\label{sec:results}

Present the findings. Refer to every table and figure in the text, for
example Table~\ref{tab:results} and Figure~\ref{fig:example}.

\begin{table}[ht]
  \centering
  \caption{Table captions go above the table. Use horizontal rules only.}
  \label{tab:results}
  \begin{tabular}{lccc}
    \toprule
    Method        & Precision & Recall & F1 \\
    \midrule
    Baseline      & 0.81      & 0.77   & 0.79 \\
    Proposed      & 0.88      & 0.85   & 0.86 \\
    \bottomrule
  \end{tabular}
\end{table}

\begin{figure}[ht]
  \centering
  \includegraphics[width=0.6\linewidth]{example-image}
  \caption{Figure captions go below the figure. Supply figures at 300\,dpi or
  as vector graphics, and make sure text in them is legible at print size.}
  \label{fig:example}
\end{figure}

\section{Discussion}
\label{sec:discussion}

Interpret the results, compare them with earlier work, and state the
limitations of the study and its practical implications.

\section{Conclusion}
\label{sec:conclusion}

Summarise the contribution and suggest directions for future research.

\begin{declarations}
  \declaration{Funding} Name the funders and grant numbers, or state that the
  research received no external funding. Anonymise funder details that would
  identify the authors during review.
  \declaration{Conflicts of interest} Declare any competing interests, or
  state that there are none.
  \declaration{Data availability} State where the data supporting the
  findings can be found, or why they cannot be shared.
  \declaration{Ethics approval} Give the approving body and reference for
  research involving human participants or animals, or state that approval
  was not required.
  \declaration{Use of generative AI} Describe any substantive use of
  generative AI tools in the research or writing, or state that none was used.
  \declaration{Author contributions} Add these in the title page for review;
  they are moved here in the accepted version.
\end{declarations}

\bibliography{references}

\end{document}
